\chapter{量子接口的端到端仿真模型与实现}[End-to-End Simulation Model and Implementation for Quantum Interfaces]

本章给出与当前模型路径一一对应的物理描述与数值流程。目标是建立从第一性原理出发、到可计算门操作结束的闭环：同一组参数既进入哈密顿量与信道，也进入最终可观测指标（成功率、真假成功分解、条件保真度、CHSH）。本章叙述严格对应当前主路径：发射态端演化 + 显式 QFC/滤波记忆链 + 光纤/BSM 在 Heisenberg 端推回 POVM。

\section{从 QED 到四能级原子-腔有效模型}[From QED to the Four-Level Atom-Cavity Effective Model]

以量子电动力学拉氏量为起点：
\begin{equation}
\mathcal{L}=\bar{\psi}\left(i\gamma^\mu\partial_\mu-m\right)\psi-eA_\mu\bar{\psi}\gamma^\mu\psi-\frac{1}{4}F_{\mu\nu}F^{\mu\nu}.
\end{equation}
在中性原子量子接口的能量与尺度范围内，经过低能近似、绝热消元、长波偶极近似、有限能级截断与旋转波近似，可得到“原子多能级 + 腔模 + 驱动 + 损耗”的有效模型。若先在二能级极限下保留一条原子跃迁与一条腔模，则退化为 Jaynes-Cummings 哈密顿量 \cite{JaynesCummings1963}：
\begin{equation}
H_{\mathrm{JC}}=\frac{\omega_{ab}}{2}\sigma_z+\omega_c\left(a^\dagger a+\frac{1}{2}\right)+g\left(\sigma^+a+\sigma^-a^\dagger\right).
\end{equation}

本项目采用通信原子的四能级有效基序
\begin{equation}
\{|0\rangle,|1\rangle,|e\rangle,|u\rangle\},
\end{equation}
其中 $|0\rangle,|1\rangle$ 为计算基，$|u\rangle$ 为受控激发前的辅助基态，$|e\rangle$ 为光学激发态。本文原子基顺序固定为 $(|0\rangle,|1\rangle,|e\rangle,|u\rangle)$。驱动仅耦合 $|u\rangle\leftrightarrow|e\rangle$，原子哈密顿量写为
\begin{equation}
H_{\mathrm{atom}}(t)=
\begin{bmatrix}
0&0&0&0\\
0&0&0&0\\
0&0&\Delta_e&\Omega(t)\\
0&0&\Omega^*(t)&\Delta_u
\end{bmatrix}.
\label{eq:h-atom-4level}
\end{equation}

偏振选择定则由两个跃迁算符表示：
\begin{equation}
S_+=|0\rangle\langle e|,\qquad S_-=|1\rangle\langle e|.
\end{equation}
将 Clebsch-Gordan 系数和偏振映射吸收入矩阵 $\boldsymbol{\alpha}\in\mathbb{C}^{2\times 2}$，定义
\begin{equation}
\sigma_H=\alpha_{H,+}S_+ + \alpha_{H,-}S_-,
\qquad
\sigma_V=\alpha_{V,+}S_+ + \alpha_{V,-}S_- .
\label{eq:sigma-hv}
\end{equation}
这一步把“物理跃迁系数”与“数值门算符”直接对接，避免后续通过经验缩放调参。

腔体采用三维截断基 $\{|{\rm vac}\rangle,|H\rangle,|V\rangle\}$，其湮灭算符为
\begin{equation}
a_H=
\begin{bmatrix}
0&1&0\\
0&0&0\\
0&0&0
\end{bmatrix},
\qquad
a_V=
\begin{bmatrix}
0&0&1\\
0&0&0\\
0&0&0
\end{bmatrix}.
\label{eq:cavity-annihilation-3d}
\end{equation}
于是 12 维发射体（atom4D $\otimes$ cavity3D）哈密顿量为
\begin{align}
H_{\mathrm{em}}(t)=&
H_{\mathrm{atom}}(t)\otimes I_3 + I_4\otimes \mathrm{diag}(0,\delta_{c,H},\delta_{c,V})\notag\\
&+g\Bigl(
\sigma_H\otimes a_H^\dagger+\sigma_H^\dagger\otimes a_H
+\sigma_V\otimes a_V^\dagger+\sigma_V^\dagger\otimes a_V
\Bigr).
\label{eq:emitter-hamiltonian-12d}
\end{align}
可见输出通道与不可见塌缩通道分别写成
\begin{equation}
L_{{\rm out},H}=\sqrt{\kappa_{{\rm ex},H}}\,I_4\otimes a_H,\quad
L_{{\rm out},V}=\sqrt{\kappa_{{\rm ex},V}}\,I_4\otimes a_V,
\end{equation}
\begin{equation}
\mathcal{C}=\left\{
\sqrt{\kappa_{{\rm in},H}}\,I_4\otimes a_H,\;
\sqrt{\kappa_{{\rm in},V}}\,I_4\otimes a_V,\;
\sqrt{\gamma_+}\,S_+\otimes I_3,\;
\sqrt{\gamma_-}\,S_-\otimes I_3
\right\}.
\end{equation}

\section{连续输出场到时间 bin 发射门}[From Continuous Output Field to Time-Bin Emission Gates]

接口输出场满足输入输出框架 \cite{GardinerCollett1985}。对每个偏振 $\mu\in\{H,V\}$，定义连续场算符 $b_\mu(t)$ 并离散为
\begin{equation}
b_{\mu,n}=\frac{1}{\sqrt{\Delta t}}\int_{t_n}^{t_n+\Delta t}b_\mu(t)\,dt,\qquad
[b_{\mu,m},b_{\nu,n}^\dagger]=\delta_{\mu\nu}\delta_{mn}.
\end{equation}

在短时步长 $\Delta t$ 下，发射门可写为
\begin{equation}
U_{\mathrm{emit}}^{(n)}
=\exp\left[
\sqrt{\Delta t}\sum_{\mu=H,V}
\left(L_{{\rm out},\mu}\otimes b_{\mu,n}^\dagger
-L_{{\rm out},\mu}^\dagger\otimes b_{\mu,n}\right)
-i\Delta t\,H_{\mathrm{em}}(t_n)\otimes I
\right].
\label{eq:emit-gate-continuous}
\end{equation}

工程实现中，每个时间站点采用 5 维 bin 基序
\begin{equation}
\{|{\rm vac}\rangle,\;|H_{780}\rangle,\;|V_{780}\rangle,\;|H_{1517}\rangle,\;|V_{1517}\rangle\}.
\end{equation}
因此 \eqref{eq:emit-gate-continuous} 先在 $(\mathrm{emitter12D}\otimes 780\mathrm{-subspace3D})$ 构造，再嵌入 bin 的 5D 空间，1517 子块保持单位，保证发射阶段只向 780 通道注入光子振幅。

为兼顾真实性与计算成本，本工作把“可见发射”保持在态端演化，而把后续线性光学与损耗映射到测量端；发射阶段仍保留不可见塌缩通道（内损与自由空间泄漏）以反映真实发射概率。这一分工与 time-bin MPS 的碰撞模型表述一致 \cite{Pichler2016TimeDelaysMPS,Wang2017MPSCascadedQuantumSystems}。

\section{QFC 与滤波记忆核的显式张量门}[Explicit Tensor Gates for QFC and Filter Memory Kernel]

\paragraph{QFC 的 5D 局域门}

QFC 在强泵近似下可写成“同偏振两频率模式间的分束器型耦合” \cite{Ikuta2018PolarizationInsensitiveQFC}。在 5D bin 基序下，$U_{\mathrm{QFC}}$ 仅作用于 $(|H_{780}\rangle,|H_{1517}\rangle)$ 与 $(|V_{780}\rangle,|V_{1517}\rangle)$ 两个二维子块：
\begin{equation}
U_{\mathrm{QFC}}=
\begin{bmatrix}
1 & 0 & 0 & 0 & 0\\
0 & c_H & 0 & -e^{-i\phi_H}s_H & 0\\
0 & 0 & c_V & 0 & -e^{-i\phi_V}s_V\\
0 & e^{i\phi_H}s_H & 0 & c_H & 0\\
0 & 0 & e^{i\phi_V}s_V & 0 & c_V
\end{bmatrix},
\end{equation}
其中 $c_{\mu}=\cos\theta_\mu,\; s_{\mu}=\sin\theta_\mu$，转换效率满足 $\eta_\mu=\sin^2\theta_\mu$。

\paragraph{滤波腔记忆模与 5D$\times$3D 步进门}

参照波导 QED 的离散记忆表示 \cite{ArranzRegidor2021WaveguideQED}，先写单偏振滤波腔的输入输出方程
\begin{equation}
\dot a_f(t)= -\left(\frac{\kappa_f}{2}+i\delta\right)a_f(t)+\sqrt{\kappa_f}\,b_{\mathrm{in}}(t),\qquad
b_{\mathrm{out}}(t)=b_{\mathrm{in}}(t)-\sqrt{\kappa_f}\,a_f(t).
\label{eq:filter-io-continuous}
\end{equation}
在时间 bin 离散后，令 $\Delta\nu=\kappa_f/(2\pi)$，可写成
\begin{equation}
a_{n+1}=r\,a_n+t\,b_n,\qquad
\tilde b_n=r^*\,b_n-t\,a_n,
\label{eq:filter-recursion-discrete}
\end{equation}
其中
\begin{equation}
r=\exp\!\left[-\left(\pi\Delta\nu+i2\pi\delta\right)\Delta t\right],\qquad
t=\sqrt{1-|r|^2},
\label{eq:filter-rt}
\end{equation}
满足 $|r|^2+t^2=1$。

为显式携带跨 bin 关联，A/B 两臂各引入一个 3D 记忆模
\begin{equation}
\{|{\rm vac}\rangle,|H\rangle,|V\rangle\}_{\rm mem},
\end{equation}
并在每一步对 $(\mathrm{bin5D}\otimes\mathrm{mem3D})$ 施加 $15\times 15$ 幺正门。式 \eqref{eq:filter-recursion-discrete} 对应的最小幺正扩展只在两个偏振子块非平凡：
\begin{equation}
\{|H_{1517},{\rm vac}_{\rm mem}\rangle,\;|{\rm vac}_{\rm bin},H_{\rm mem}\rangle\}
\end{equation}
和
\begin{equation}
\{|V_{1517},{\rm vac}_{\rm mem}\rangle,\;|{\rm vac}_{\rm bin},V_{\rm mem}\rangle\},
\end{equation}
二者均为
\begin{equation}
\begin{bmatrix}
r&-t\\
t&r^*
\end{bmatrix}.
\label{eq:filter-step-subblock}
\end{equation}
从张量角度，步进门写为
\begin{equation}
U_{\mathrm{step}}
=I_{15}
\oplus
\left(
\begin{bmatrix}
r&-t\\
t&r^*
\end{bmatrix}_{H}
\oplus
\begin{bmatrix}
r&-t\\
t&r^*
\end{bmatrix}_{V}
\right),
\label{eq:filter-step-tensor}
\end{equation}
因此“当前 bin 的 1517 分量”和“记忆模占据”在每步进行相干交换并积累跨 bin 相关，形成非马尔可夫尾迹 \cite{ArranzRegidor2021WaveguideQED,Garraway1997DecayStrongReservoir}。

\paragraph{链布局与唯一路径实现}

本工作固定采用显式记忆路径。发射结束布局为
\begin{equation}
{\rm atomA,\ atomB,\ A1,\ B1,\ldots, AN,\ BN}.
\end{equation}
QFC+滤波阶段引入两臂记忆模，逐 bin 相邻作用后收拢为唯一后续布局
\begin{equation}
{\rm atomA,\ atomB,\ memA,\ memB,\ A1,\ B1,\ldots, AN,\ BN}.
\label{eq:final-chain-layout}
\end{equation}
该布局直接服务后续收缩：原子与记忆核在链首，时间 bin 区域结构规则，降低索引管理与 SWAP 调度复杂度。

\section{光纤与 BSM 的 Heisenberg-POVM 统一建模}[Unified Heisenberg-POVM Modeling for Fiber and BSM]

\paragraph{探测端 effect 构造与可分辨性混合}

对每个偏振 $\mu\in\{H,V\}$，50:50 分束器先定义输入输出模变换
\begin{equation}
\begin{bmatrix}
c_{1\mu}\\
c_{2\mu}
\end{bmatrix}
=
\frac{1}{\sqrt{2}}
\begin{bmatrix}
1&i\\
i&1
\end{bmatrix}
\begin{bmatrix}
a_{A\mu}\\
a_{B\mu}
\end{bmatrix}.
\label{eq:bs-mode-map}
\end{equation}
在探测端口上，对单个探测器 $j$（效率 $\eta_j$、暗计数概率 $d_j$）可写二元 effect
\begin{equation}
E_j^{(1)}=I-(1-d_j)\left(I-\eta_j n_j\right),\qquad
E_j^{(0)}=I-E_j^{(1)}.
\label{eq:single-detector-effect}
\end{equation}
四路联合记录 $\mathbf{s}=(s_1,s_2,s_3,s_4)$ 的 effect 为
\begin{equation}
E_{\mathbf{s}}=\bigotimes_{j=1}^{4}E_j^{(s_j)}.
\label{eq:joint-detection-effect}
\end{equation}
BSM 记录集由式 \eqref{eq:joint-detection-effect} 线性组合得到，随后通过分束器对偶映射
\begin{equation}
E\leftarrow U_{\mathrm{BS}}^\dagger E U_{\mathrm{BS}}
\end{equation}
推回输入端。为描述部分可分辨性，采用“不可分辨 3D 映射”和“可分辨 9D 标签映射”的凸混合：
\begin{equation}
E^{(v_{\mathrm{res}})}=v_{\mathrm{res}}E_{\mathrm{int}}+(1-v_{\mathrm{res}})E_{\mathrm{dist}},
\label{eq:vres-mixture}
\end{equation}
其中 $v_{\mathrm{res}}\in[0,1]$。

\paragraph{光纤与相位噪声的对偶推回}

光纤段包含偏振旋转（Jones）、偏振相关损耗（PDL）与相位剖面。每个 bin 的相位写为
\begin{equation}
\phi_n=\phi_{\mathrm{slope}}\!\left(n-\frac{N-1}{2}\right)+\xi_n,\qquad \xi_n\sim\mathcal{N}(0,\sigma_\phi^2),
\end{equation}
并通过 $U_B(\phi_n)$ 进入 B 臂 Jones 矩阵。所有光纤和探测信道均通过对偶映射
\begin{equation}
\Lambda^\dagger(E)=\sum_k K_k^\dagger E K_k
\label{eq:channel-adjoint}
\end{equation}
作用于 effect，而非作用在态上。这样可在保持纯态 MPS 主态演化的同时，严格保留损耗与噪声统计。

\paragraph{成功判据与核心指标}

设 $\mathcal{S}$ 为窗口内 BSM 成功模式集合。枚举阶段直接从收缩获得
\begin{align}
p_s &= \sum_{m\in\mathcal{S}} \mathrm{Tr}(E_m\rho),\\
p_t &= \sum_{m\in\mathcal{S}} \mathrm{Tr}(E_m^{\mathrm{true}}\rho),\\
p_f &= p_s-p_t,
\end{align}
并输出
\begin{equation}
p_{t|11}=\frac{p_t}{p_{11}}.
\end{equation}
该指标体系同时给出“宣告成功率”“真成功分量”与“假成功污染”，可直接用于误差预算。

\section{实现口径、复杂度与可扩展性}[Conventions, Complexity, and Scalability]

\paragraph{统一口径与复杂度瓶颈}

当前实现对时间口径做了显式统一：哈密顿量参数（如 $\Omega,g,\Delta$）按 $\mathrm{rad/s}$ 解释，耗散速率（如 $\kappa,\gamma$）按 $\mathrm{1/s}$ 解释；在进入发射内核前统一换算。该约束用于避免 $2\pi$ 口径混用导致的时间尺度偏差，并已用于发射波包标定（目标时间常数接近实验 $26.2\,\mathrm{ns}$ \cite{vanleent2022entangling}）。
为了使这一定义在不同离散步长、不同窗口大小和不同参数扫描任务中保持可比，本工作同时约束了“时间积分相关无量纲组合”的一致性：当 $\Delta t$ 改变时，影响动力学的关键组合量（例如 $g\sqrt{\Delta t}$ 与 $\kappa\Delta t$）必须与连续模型保持对应，不能让数值网格本身重定义物理时间尺度。也正因此，本章所有运行时间解释都以“固定物理口径”作为前提，而非以“固定步数”作为前提。

在当前参数规模（$N\!=\!100$、window $\sim 70$ ns）下，主耗时通常集中在 POVM 枚举收缩；发射和 QFC 阶段开销次之。总复杂度可粗略写为
\begin{equation}
T_w\approx T_e+T_q+T_n+T_m,
\end{equation}
其中
\begin{equation}
T_n\propto O\!\left(N\cdot W\cdot C(\chi)\right),
\end{equation}
$W$ 为窗口宽度对应的 bin 数，$C(\chi)$ 表示给定键维下的单次收缩成本。该结构解释了“参数几乎不变但窗口或枚举口径改变时墙钟显著变化”的现象。进一步地，若把状态更新、effect 推回与最终枚举看成三层管线，则 wall-clock 的主导项在不同任务类型下会迁移：参数扫描任务通常受 $T_n$ 主导，而单次高精度结构分析则更受 $C(\chi)$ 的局域高秩峰值主导。这一分解直接指导第四章中的性能诊断与优化优先级设定。

\section{本章小结}[Chapter Summary]

本章从 QED 出发，建立了端到端模型：四能级原子-腔发射门、5D QFC 门、5D$\times$3D 滤波记忆门、以及 Heisenberg 端的光纤与 BSM-POVM 收缩。核心技术路线是“态端保留发射与记忆核，线性光学和损耗推回 effect”。其直接收益是：在不牺牲物理可解释性的前提下，将复杂链路压缩为可枚举、可分解、可做误差预算的计算流程，为下一章的实验对照结果与参数扫描提供统一基线。
